% ------------------------------------------------------------
\escenario{ESC-18}{La plataforma de distribución revisa el paquete}

\begin{tabla}{|L{3.2cm}|Y|}
\fichacab
\bloque{Trazabilidad}
\parte{Característica}{QA-10 Flexibilidad: instalabilidad (condicionada a Q-10).}
\parte{Stakeholders}{STK-19 Plataformas de distribución, STK-14 Equipo de operaciones y despliegue.}
\parte{Concerns}{CRN-25: \emph{``Antes de aceptar el juego en nuestro catálogo exigimos clasificación por edad, un empaquetado conforme a nuestras normas y una política de chat que cumpla la regulación de cada región.''}}
\parte{Restricciones y dependencias}{CON-16 (plataforma .NET y su runtime), DEP-08 (plataformas de distribución).}
\parte{Tensiones y preguntas}{Q-10 (publicación en una plataforma de distribución). Solo aplica si se decide publicar.}
\parte{Tipo y prioridad}{De uso, en tiempo de instalación. Prioridad (B, B).}
\bloque{Escenario}
\parte{Fuente del estímulo}{Plataforma de distribución.}
\parte{Estímulo}{Recibe el paquete del juego y lo instala en sus equipos de prueba para decidir si lo acepta en su catálogo.}
\parte{Ambiente}{Equipos de prueba limpios de la plataforma, sin el runtime de .NET instalado de antemano.}
\parte{Artefacto}{Paquete de instalación del cliente de juego.}
\parte{Respuesta}{El juego se instala, se abre y se desinstala sin pedir componentes adicionales, y el paquete declara lo que necesita tal como lo piden las normas de la plataforma, incluida la clasificación por edad.}
\parte{Medida de la respuesta}{El paquete se aprueba en la primera revisión sin observaciones de empaquetado. En el 100~\% de los equipos de prueba el juego se instala y se abre sin descargar nada aparte, y al desinstalarlo no deja archivos.}
\end{tabla}

\textbf{Por qué es relevante.} Si se decide publicar (Q-10), la plataforma tiene la última palabra: un paquete rechazado significa que el juego no llega a sus usuarios, sin importar lo bien que funcione. El riesgo concreto viene de CON-16: un juego en .NET que da por hecho que el runtime ya está instalado falla en un equipo limpio, que es justo como prueba la plataforma. Su prioridad es baja porque depende de una pregunta abierta, pero se deja formulado para que, si Q-10 se resuelve a favor, el criterio de aceptación ya exista.

\textbf{Cómo se verifica.} Se prepara una máquina virtual limpia con el sistema operativo objetivo y sin .NET, se instala el paquete, se abre el juego y se juega una partida, se desinstala y se compara la lista de archivos del equipo antes y después. Se revisa el paquete contra la lista de requisitos publicada por la plataforma.
